\unit{Exposición 232 -- Esquemas de Picard: teoremas de existencia}
\sid{232}
\src{143}
\seminares{14.\textsuperscript{o} año, 1961/62, n.\textsuperscript{o} 232\\Febrero de 1962}
\paperhead{TÉCNICA DE DESCENSO Y TEOREMAS DE EXISTENCIA EN GEOMETRÍA ALGEBRAICA\\
V. ESQUEMAS DE PICARD: TEOREMAS DE EXISTENCIA}%
  {por Alexander GROTHENDIECK}

\fsection{1}{Grupos y funtores de Picard relativos}

Para todo preesquema (más generalmente, todo espacio anillado) $X$, llamamos al
grupo de clases, consideradas salvo isomorfismo, de Módulos invertibles (es
decir, Módulos localmente isomorfos a $\mathcal O_X$) sobre $X$ el
\emph{grupo de Picard} (absoluto) de $X$, y lo denotamos por
$\operatorname{Pic}(X)$. Se tiene, por tanto, un isomorfismo canónico
\begin{equation}
\operatorname{Pic}(X)\xrightarrow{\sim}H^1(X,\mathcal O_X^*),
\tag{1.1}
\end{equation}
donde $\mathcal O_X^*$ denota el haz de unidades de $\mathcal O_X$ (que se
identifica, en efecto, con el haz de automorfismos del Módulo invertible
estándar $\mathcal O_X$). Es claro que
$X\longmapsto\operatorname{Pic}(X)$ es un funtor contravariante en $X$, y que
el isomorfismo (1.1) es funtorial.

Si $X$ es un preesquema sobre otro preesquema $S$, al variar $S'$ en la
categoría $(\mathrm{Sch})/S$ de preesquemas sobre $S$, la construcción anterior
proporciona un funtor contravariante
$S'\longmapsto\operatorname{Pic}(X\times_SS')$. Este funtor no tiene ninguna
posibilidad de ser «representable» ([4], II, A), pues los automorfismos no
triviales de los Módulos invertibles que se pretende clasificar impiden que sea
de «naturaleza local» ([5], IV, 5.4). Hay, pues, que «hacerlo local». De modo
general, para todo preesquema relativo $X/S$, introduzcamos el grupo relativo
\begin{equation}
\operatorname{Pic}'(X/S)=H^0\bigl(S,R^1f_*(\mathcal O_X^*)\bigr),
\tag{1.2}
\end{equation}
donde $f:X\longrightarrow S$ es el morfismo estructural (compárese [4], II,
C 3). En el pasaje citado se llama a este grupo grupo de Picard relativo; por
razones que aparecerán enseguida, conviene llamarlo aquí \emph{grupo de Picard
relativo restringido} de $X/S$. Al variar $S'$ en $(\mathrm{Sch})/S$, la
asignación $S'\longmapsto\operatorname{Pic}'(X\times_SS'/S')$ es un funtor
contravariante en $S'$, denotado asimismo por $\operatorname{Pic}'_{X/S}$ y
definido, por tanto, esencialmente por
\begin{equation}
\operatorname{Pic}'_{X/S}(S')=
\operatorname{Pic}'(X\times_SS'/S').
\tag{1.3}
\end{equation}

\src{144}
Este funtor es ahora de «naturaleza local», pues se ha efectuado el ajuste
necesario. Intuitivamente, el segundo miembro de (1.3) se interpreta como el
conjunto de «familias algebraicas» de clases de haces invertibles sobre (las
fibras de) $X/S$, indexadas por el preesquema de parámetros $S'/S$. Cuando el
funtor $\operatorname{Pic}'$ es representable, el preesquema sobre $S$ que lo
representa se denota por $\operatorname{Pic}'_{X/S}$ y se llama
\emph{preesquema de Picard} de $X$ sobre $S$; se tiene entonces
\begin{equation}
\operatorname{Hom}_S(S',\operatorname{Pic}'_{X/S})
\simeq\operatorname{Pic}'_{X/S}(S')
=\operatorname{Pic}'(X\times_SS'/S').
\tag{1.4}
\end{equation}

Hay, sin embargo, casos importantes en los que $\operatorname{Pic}'_{X/S}$ no
es representable (por ejemplo, una variedad de «Brauer--Severi» sobre un
cuerpo $k$ sin punto racional sobre $k$), aunque sigue habiendo una definición
natural de preesquema de Picard relativo. Ello se debe a que, al definir el
funtor $\operatorname{Pic}'$ a partir de los grupos de Picard absolutos
$\operatorname{Pic}(X\times_SS'/S')$, aún no se ha localizado suficientemente;
más precisamente, $\operatorname{Pic}'$ no es, en general, «compatible con el
descenso fielmente plano». Hagámoslo explícito.

Sea $\mathfrak M$ la clase de los morfismos de preesquemas que son
\emph{fielmente planos y cuasicompactos}; esta clase es estable por cambio de
base y por composición. Sea $P$ un funtor contravariante de $(\mathrm{Sch})/S$
a la categoría de conjuntos y, para todo $S$-morfismo
$u:T'\longrightarrow T$ con $u\in\mathfrak M$, consideremos el diagrama
\begin{equation}
P(T)\longrightarrow P(T')\rightrightarrows P(T'\times_TT')
\tag{1.5}
\end{equation}
obtenido al aplicar $P$ a
\[
T\longleftarrow T'\mathrel{\substack{\xleftarrow{\,\mathrm{pr}_1\,}\\[-0.35ex]
\xleftarrow[\,\mathrm{pr}_2\,]{}}}T'\times_TT'.
\]
Si $P$ es representable, la teoría del descenso ([4], I, B, teorema 2) implica
que (1.5) es exacto para todo $u\in\mathfrak M$. Esto se expresa diciendo que
$P$ es compatible con $\mathfrak M$; en este caso, que $P$ es «compatible con
el descenso fielmente plano», o también que el «prehaz» $P$ sobre
$(\mathrm{Sch})/S$ es un «haz» para la noción de localización proporcionada
por $\mathfrak M$. Para un $P$ arbitrario, un procedimiento estándar, familiar
por la localización topológica ordinaria, permite asociarle un «haz»
$\mathcal P$ y un homomorfismo de funtores $P\longrightarrow\mathcal P$
universal en el sentido evidente. El cálculo de $\mathcal P$ puede hacerse
explícito como sigue. Para definir $\mathcal P(T)$, para cada $T'$ sobre $T$
tal que $u:T'\longrightarrow T$ pertenezca a $\mathfrak M$, sea

\src{145}
$\overline H^0(T'/T,P)$ el subconjunto de $P(T')$ formado por los elementos
$\xi$ cuyas dos imágenes $\xi_1,\xi_2$ en $P(T'\times_TT')$ tienen la misma
imagen en $P(T'')$ para algún morfismo
$v:T''\longrightarrow T'\times_TT'$ con $v\in\mathfrak M$. [N.B. -- El
conjunto $\overline H^0$ así definido es, por tanto, mayor que el conjunto
$H^0(T'/T,P)$ introducido en [4], I, A, 4(a).] Al variar $T'$ sobre un $T$ fijo
(siempre con $u\in\mathfrak M$), los $\overline H^0(T'/T,P)$ forman un sistema
inductivo cuando se dota al conjunto de los $T'$ del preorden definido por la
dominación, y se pone
\begin{equation}
\mathcal P(T)=\varinjlim_{T'}\overline H^0(T'/T,P).
\tag{1.6}
\end{equation}
La ley funtorial en $T$ de esta expresión es inmediata.

Cuando
\[
P(T)=\operatorname{Pic}(X\times_ST),
\]
el funtor contravariante sobre $(\mathrm{Sch})/S$ definido por (1.6) se llama
\emph{funtor de Picard relativo} de $X$ sobre $S$ y se denota por
$\operatorname{Pic}_{X/S}$; el grupo $\operatorname{Pic}_{X/S}(S)$ se llama
\emph{grupo de Picard relativo} de $X$ sobre $S$ y se denota por
$\operatorname{Pic}(X/S)$. Se obtiene entonces la biyección evidente
\begin{equation}
\operatorname{Pic}_{X/S}(T)\simeq\operatorname{Pic}(X\times_ST/T).
\tag{1.7}
\end{equation}

Así, un elemento de $\operatorname{Pic}(X/S)$ viene definido por un elemento
$\xi'$ de un grupo $\operatorname{Pic}(X\times_SS')$, donde
$S'\longrightarrow S$ es fielmente plano y cuasicompacto, tal que existe un
morfismo fielmente plano y cuasicompacto
$S''\longrightarrow S'\times_SS'$ para el cual coinciden las dos imágenes
inversas de $\xi'$ en $\operatorname{Pic}(X\times_SS'')$. Un elemento $\xi'$ de
$\operatorname{Pic}(X\times_SS')$ y un elemento $\xi_1$ de
$\operatorname{Pic}(X\times_SS_1)$ que satisfagan las condiciones anteriores
definen el mismo elemento de $\operatorname{Pic}(X/S)$ si y solo si existe un
morfismo fielmente plano y cuasicompacto
$S'_1\longrightarrow S'\times_SS_1$ tal que las imágenes de ambos elementos en
$\operatorname{Pic}(X\times_SS'_1)$ sean iguales. A menudo sigue siendo
conveniente trabajar con el funtor $P'=\operatorname{Pic}'_{X/S}$ introducido
antes; se comprueba inmediatamente que el morfismo canónico
$P\longrightarrow P'$ induce un isomorfismo
\begin{equation}
\mathcal P\xrightarrow{\sim}\mathcal P',
\tag{1.8}
\end{equation}
lo que proporciona una descripción, con frecuencia más cómoda, de
$\operatorname{Pic}_{X/S}$ en términos de
$\operatorname{Pic}'_{X/S}=P'$. En virtud de 2.3, al sustituir $P$ por $P'$ en
la descripción anterior de $\operatorname{Pic}(X/S)$, se puede de hecho
\src{146}
tomar $S''=S'\times_SS'$ y $S'_1=S'\times_SS_1$, al menos bajo las condiciones
especificadas en el pasaje citado.

Cuando el funtor $\operatorname{Pic}_{X/S}$ es representable, se dice que
$X/S$ admite un preesquema de Picard, y el preesquema sobre $S$ que representa
el funtor se llama \emph{preesquema de Picard} de $X$ sobre $S$ y sigue
denotándose por $\operatorname{Pic}_{X/S}$. Basta evidentemente que
$P'=\operatorname{Pic}'_{X/S}$ sea representable, pues entonces $P'$ ya es un
«haz», y (1.8) muestra que $P'\longrightarrow\mathcal P'$ es el morfismo
canónico
\begin{equation}
\operatorname{Pic}'_{X/S}\longrightarrow\operatorname{Pic}_{X/S},
\tag{1.9}
\end{equation}
que es entonces un isomorfismo. Nuestra terminología concuerda así con la
introducida anteriormente en (1.4). En general, cuando
$\operatorname{Pic}_{X/S}$ existe, queda definido por el isomorfismo funtorial
\begin{equation}
\operatorname{Hom}_S(S',\operatorname{Pic}_{X/S})\xrightarrow{\sim}
\operatorname{Pic}(X\times_SS'/S').
\tag{1.10}
\end{equation}

\fsection{2}{Relaciones entre los diversos grupos de Picard relativos y absolutos}

\result{PROPOSICIÓN}{2.1} Sea $f:X\longrightarrow S$ un morfismo tal que
$\mathcal O_S\xrightarrow{\sim}f_*(\mathcal O_X)$. Se tiene entonces una
sucesión exacta
\[
0\longrightarrow\operatorname{Pic}(S)\longrightarrow\operatorname{Pic}(X)
\longrightarrow\operatorname{Pic}'(X/S).
\]
Cuando $X$ admite una sección sobre $S$, el último morfismo es sobreyectivo;
es decir, se tiene un isomorfismo
\[
\operatorname{Pic}'(X/S)\xrightarrow{\sim}
\operatorname{Pic}(X)/\operatorname{Pic}(S).
\]

La sucesión exacta puede considerarse como la sucesión exacta en grados bajos
asociada a la sucesión espectral de Leray para $f$ y $\mathcal O_X^*$. La
segunda afirmación es igualmente formal.

\result{PROPOSICIÓN}{2.2} Sea $f:X\longrightarrow S$ un morfismo cuasicompacto
y separado tal que $\mathcal O_S\xrightarrow{\sim}f_*(\mathcal O_X)$, y sea
$S'\longrightarrow S$ fielmente plano y cuasicompacto. Entonces:

\noindent (i) $\operatorname{Pic}'(X/S)\longrightarrow
\operatorname{Pic}'(X\times_SS'/S')$ es inyectivo;

\noindent (ii) si $X$ admite localmente una sección sobre $S$ (es decir, todo
$s\in S$ posee un entorno abierto $U$ tal que $X|U$ admite una sección sobre
$U$), entonces el diagrama
\[
\operatorname{Pic}'(X/S)\longrightarrow
\operatorname{Pic}'(X\times_SS'/S')\rightrightarrows
\operatorname{Pic}'(X\times_SS''/S'')
\]

\src{147}
donde $S''=S'\times_SS'$, es exacto.

La primera afirmación se deduce de las propiedades elementales del descenso
fielmente plano y de la observación general siguiente. Si
$f:X\longrightarrow S$ satisface
$\mathcal O_S\xrightarrow{\sim}f_*(\mathcal O_X)$, entonces el funtor
$\mathcal F\longmapsto f^*(\mathcal F)$ de los Módulos localmente libres de
tipo finito sobre $S$ a los Módulos localmente libres de tipo finito sobre $X$
es plenamente fiel, y su imagen esencial está formada por los Módulos
$\mathcal G$ sobre $X$ tales que $f_*(\mathcal G)$ es localmente libre y el
homomorfismo canónico
\[
f^*\bigl(f_*(\mathcal G)\bigr)\longrightarrow\mathcal G
\]
es un isomorfismo. La segunda afirmación fue demostrada mediante la teoría del
descenso en [4], I, B, 4.

Los resultados de 2.2 también pueden enunciarse como sigue.

\result{COROLARIO}{2.3} Bajo las condiciones de 2.2, el homomorfismo canónico
(1.9) $\operatorname{Pic}'_{X/S}\longrightarrow\operatorname{Pic}_{X/S}$ es
inyectivo, y es biyectivo si $X$ admite localmente una sección sobre $S$.
(Así, en este último caso, el grupo de Picard relativo
$\operatorname{Pic}(X/S)$ y el grupo de Picard relativo restringido
$\operatorname{Pic}'(X/S)$ son idénticos.)

Combinando esto con 2.1 se obtiene:

\result{COROLARIO}{2.4} Bajo las condiciones de 2.2, se tiene una sucesión
exacta
\[
0\longrightarrow\operatorname{Pic}(S)\longrightarrow\operatorname{Pic}(X)
\longrightarrow\operatorname{Pic}(X/S).
\]
Cuando $X$ admite una sección sobre $S$, el último homomorfismo es
sobreyectivo; es decir, se tiene entonces un isomorfismo
\[
\operatorname{Pic}(X/S)\xrightarrow{\sim}
\operatorname{Pic}(X)/\operatorname{Pic}(S).
\]

\result{OBSERVACIÓN}{2.5} Sea $f:X\longrightarrow S$ un morfismo tal que
$\mathcal O_S\xrightarrow{\sim}f_*(\mathcal O_X)$, y sea $g$ una sección de
$X$ sobre $S$. Para un Módulo invertible $\mathcal L$ sobre $X$, llamemos al
isomorfismo $\mathcal O_S\xrightarrow{\sim}g^*(\mathcal L)$ una
$g$-rigidificación de $\mathcal L$, y llamemos a un Módulo invertible
$\mathcal L$ sobre $X$ provisto de una $g$-rigidificación \emph{Módulo
invertible $g$-rigidificado}. Todo automorfismo de una estructura de este
tipo es trivial, y $\operatorname{Pic}'(X/S)$ se identifica con el conjunto de
clases de isomorfismo de Módulos invertibles $g$-rigidificados sobre
$X$\footnote{El texto impreso dice «sobre $S$»; la rigidificación mediante una
sección $g:S\to X$ concierne aquí al Módulo $\mathcal L$ sobre $X$.}. (Este es
el hecho que permitió emplear la teoría del descenso en la demostración de
2.2(ii).) Esto da una nueva interpretación de $\operatorname{Pic}(X/S)$, al
menos cuando $f$ es además cuasicompacto y separado, de modo que
$\operatorname{Pic}(X/S)\simeq\operatorname{Pic}'(X/S)$ por 2.3.

\src{148}
\result{OBSERVACIÓN}{2.6} Sea $f:X\longrightarrow S$ como en 2.2, y sea
$S'\longrightarrow S$ fielmente plano y cuasicompacto, de modo que exista un
$S$-morfismo $S'\longrightarrow X$, o, equivalentemente, que
$X'=X\times_SS'$ admita una sección sobre $S'$. Pongamos
$S''=S'\times_SS'$ y $X''=X\times_SS''$, y consideremos la sucesión exacta
\[
\operatorname{Pic}(X/S)\longrightarrow\operatorname{Pic}(X'/S')
\rightrightarrows\operatorname{Pic}(X''/S'').
\]
Aplicando 2.4 se obtiene una sucesión exacta
\[
\operatorname{Pic}(X/S)\longrightarrow
\operatorname{Pic}(X')/\operatorname{Pic}(S')
\rightrightarrows
\operatorname{Pic}(X'')/\operatorname{Pic}(S''),
\]
y, en particular, todo elemento del grupo de Picard relativo ya «proviene»
de un elemento de $\operatorname{Pic}(X')$. Esto simplifica sustancialmente la
descripción del grupo de Picard relativo dada en la sección anterior, y también
la del funtor de Picard de $X$ sobre $S$, pues para todo $T$ sobre $S$ el
argumento precedente se aplica a $X\times_ST/T$ y al morfismo
$T'=S'\times_ST\longrightarrow T$. Si, por ejemplo, $f$ mismo es fielmente
plano, se puede tomar $S'=X$. Cuando $f$ es además de tipo finito
(respectivamente liso, etc.), esto permite limitarse, en la descripción del
funtor de Picard relativo
$T\longmapsto\operatorname{Pic}(X\times_ST/T)$, a cambios de base
$T'\longrightarrow T$ que sean de tipo finito (respectivamente lisos, etc.).
Cuando $f$ es proyectivo y plano y $S$ es localmente noetheriano, se demuestra
que el morfismo $S'\longrightarrow S$ anterior puede elegirse de modo que $S'$
sea la unión disjunta de recubrimientos planos $S'_i$ de abiertos $S_i$ que
cubren $S$; si $f$ es además separable, los $S'_i$ pueden tomarse étale sobre
$S_i$.

\fsection{3}{El teorema principal de existencia: enunciado}

No se dispone, ni siquiera a título de conjetura, de un enunciado de existencia
de preesquemas de Picard que englobe todos los casos conocidos. Una condición
«prácticamente necesaria», si puede decirse así, es que
$f:X\longrightarrow S$ sea propio, asegurando las propiedades esenciales de
finitud, y plano. Estas condiciones no son suficientes, incluso si $S$ es el
espectro del álgebra de los números duales $k[t]/(t^2)$ sobre un cuerpo $k$
(por ejemplo, el cuerpo $\mathbf C$ de los números complejos) y $X$ tiene
dimensión 1. En el momento de redactar esta exposición, los teoremas de
existencia más importantes para el preesquema de Picard se deducen del teorema
siguiente.

\result{TEOREMA}{3.1} Sea $f:X\longrightarrow S$ un morfismo de preesquemas
localmente noetherianos. Supongamos:

\noindent (i) $f$ es proyectivo;

\src{149}
\noindent (ii) $f$ es plano;

\noindent (iii) las fibras geométricas de $f$ son íntegras.

Bajo estas condiciones, $\operatorname{Pic}_{X/S}$ existe.

La demostración, que se esbozará en las dos secciones siguientes, mostrará al
mismo tiempo lo que sigue. Sea $\xi$ la sección de
$\operatorname{Pic}_{X/S}$ correspondiente a un haz muy amplio
$\mathcal O_X(1)$ sobre $X/S$, es decir, inducido por una inmersión proyectiva
$X\longrightarrow\mathbf P(\mathcal E)$. Existe entonces un abierto $U$ de
$\operatorname{Pic}_{X/S}$, unión disjunta de abiertos cuasiproyectivos sobre
$S$, tal que $U$ es estable por traslación por $\xi$, y
$\operatorname{Pic}_{X/S}$ es la unión creciente de los abiertos $U-n\xi$,
todos isomorfos a $U$. En particular, bajo las condiciones de 3.1,
$\operatorname{Pic}_{X/S}$ es separado sobre $S$.

\result{OBSERVACIÓN}{3.2} Algunos ejemplos muestran (por ejemplo, con $S$ el
espectro de un anillo de valoración discreta y $X$ de dimensión relativa 1
sobre $S$) que, si se omite la hipótesis (iii) de 3.1 y se sustituye por la
hipótesis más débil de que, para todo $s\in S$, el homomorfismo
\[
k(s)\longrightarrow H^0(X_s,\mathcal O_{X_s})
\]
es un isomorfismo, entonces $\operatorname{Pic}_{X/S}$ no es necesariamente
separado sobre $S$. Esto ocurre tanto cuando las fibras geométricas de $f$ son
reducidas pero una fibra geométrica genérica íntegra «estalla» por
especialización en dos componentes irreducibles, como cuando las fibras
geométricas de $f$ son irreducibles pero una fibra geométrica genérica íntegra
se especializa en una «fibra múltiple». El primer caso se presenta, por ejemplo,
con una cónica que degenera en dos rectas secantes; D. Mumford proporcionó un
ejemplo del segundo, con una curva elíptica que degenera en una curva elíptica
doble. Estos ejemplos son válidos en toda característica.

\result{OBSERVACIÓN}{3.3} Bajo las condiciones de 3.1, no sé si
$\operatorname{Pic}_{X/S}$ es una unión disjunta de abiertos de tipo finito, y
por tanto cuasiproyectivos, sobre $S$.\footnote{La fe de erratas de 1962 informa de
que D. Mumford respondió afirmativamente a esta cuestión mediante la teoría de
los esquemas de Picard desarrollada aquí.} Los polinomios de Hilbert
$Q\in\mathbf Q[t]$, como en el caso de los esquemas de Hilbert ([4], IV), dan
una descomposición de $\operatorname{Pic}_{X/S}$ en una unión disjunta de
abiertos $\operatorname{Pic}^{Q}_{X/S}$, y parece plausible que estos sean
cuasiproyectivos sobre $S$; esto se verá al menos en la exposición siguiente
cuando $f$ sea liso. Obsérvese que, si se sustituye la hipótesis (i) por la de
que $X$ sea localmente proyectivo sobre $S$ (lo cual basta para 3.1, pues la
existencia de $\operatorname{Pic}_{X/S}$ es manifiestamente local sobre $S$),
resulta, por el contrario, fácil

\src{150}
dar ejemplos en los que $\operatorname{Pic}_{X/S}$ contiene componentes conexas
que no son de tipo finito sobre $S$. Sea, por ejemplo, $X_0$ una variedad
algebraica no singular y proyectiva sobre un cuerpo algebraicamente cerrado
$k$, provista de un automorfismo $u$ y de un elemento $\xi$ del grupo de
Néron--Severi de $X_0$ tales que los $u^n(\xi)$ sean dos a dos distintos. Se
puede tomar para $X_0$ el producto de una curva elíptica $E$ por sí misma, y
para $u$ el automorfismo
\[
(x,y)\longmapsto(x,y+x)
\]
de $E\times E$. Sea $S$ la unión de dos curvas irreducibles no singulares que
se cortan en dos puntos $a$ y $b$. Existe un recubrimiento principal conexo de
grupo $\mathbf Z$, $P$, sobre $S$; usando las acciones de $\mathbf Z$ sobre $X_0$
definidas por $u$, se obtiene un fibrado asociado sobre $S$ de fibra $X_0$,
trivial sobre $S-a$ y $S-b$ y, en el caso particular considerado, de hecho un
esquema abeliano sobre $S$. Se comprueba fácilmente que
$\operatorname{Pic}_{X/S}$, que también es el fibrado asociado a $P$ y a las
acciones de $\mathbf Z$ sobre $\operatorname{Pic}_{X_0/k}$ por medio de $u$,
contiene una componente conexa isomorfa a
\[
P\times\operatorname{Pic}^{0}_{X_0/k},
\]
donde $\operatorname{Pic}^{0}$ denota la componente conexa de la identidad en
$\operatorname{Pic}$; esta componente no es de tipo finito sobre $S$.
[Fenómenos análogos se producen también en diversos casos de preesquemas de
Picard no separados sobre $S$, como los considerados en 3.2.]

\fsection{4}{Divisores de Cartier relativos y fibrados proyectivos}

Solo necesitaremos divisores positivos y omitiremos este calificativo adicional
en toda la sección.

Sea $X$ un preesquema. Un \emph{divisor de Cartier} sobre $X$, o simplemente un
divisor sobre $X$, es un subpreesquema cerrado $D$ de $X$ definido por un Ideal
$\mathcal J$ que es un Módulo invertible, es decir, localmente generado por una
sección que no es divisor de cero en $\mathcal O_X$. A $D$ asociamos el Módulo
invertible
\[
\mathcal L(D)=\mathcal J^{-1}.
\]
La inyección canónica $\mathcal J\longrightarrow\mathcal O_X$ proporciona un
homomorfismo canónico
\[
s_D:\mathcal O_X\longrightarrow\mathcal J^{-1}=\mathcal L(D),
\qquad\text{es decir,}\quad s_D\in\Gamma(X,\mathcal L(D)).
\]
Además, dar un divisor equivale esencialmente a dar un Módulo invertible
$\mathcal L$ sobre $X$ junto con una sección $s$ que no sea divisor de cero en
ningún punto, asociando recíprocamente a tal par $(\mathcal L,s)$ el
«divisor» de $s$, denotado por $\operatorname{div}(s)$. Para un $\mathcal L$
invertible fijo sobre $X$, el conjunto de

\src{151}
divisores $D$ que definen $\mathcal L$ está en biyección con el conjunto
cociente
\[
\Gamma(X,\mathcal L)^*/\Gamma(X,\mathcal O_X^*),
\]
donde $\Gamma(X,\mathcal L)^*$ es el subconjunto de
$\Gamma(X,\mathcal L)$ formado por las secciones que no son divisores de cero
en ningún punto.

Supongamos ahora dado un morfismo localmente de tipo finito
$f:X\longrightarrow S$, con $S$ localmente noetheriano para simplificar. Sean
$\mathcal J$ un Ideal coherente sobre $X$, $D$ el subesquema de $X$ que define,
$x\in X$, y pongamos $s=f(x)$. Las condiciones siguientes son equivalentes:

\noindent (i) $\mathcal J$ es invertible en $x$ (es decir, $\mathcal J_x$ está
generado por un elemento regular de $\mathcal O_{X,x}$), y $D$ es plano sobre
$S$ en $x$;

\noindent (ii) $X$ y $D$ son planos sobre $S$ en $x$, y $D_s$ es un divisor de
Cartier sobre la fibra $X_s$ en $x$;

\noindent (iii) $X$ es plano sobre $S$ en $x$, y $\mathcal J_x$ está generado
por un elemento $f_x$ que induce sobre $X_s$ un germen que no es divisor de
cero.

Se dice entonces que $D$ es un \emph{divisor de Cartier relativo} sobre $X/S$
(o simplemente un divisor relativo sobre $X/S$) en el punto considerado. La
condición (i) muestra que $D$ es también un divisor relativo en los puntos
vecinos de $x$. Por tanto, si $X$ y $D$ son planos sobre $S$ y $D$ es propio
sobre $S$, el conjunto de los $s\in S$ tales que $D_s$ es un divisor de Cartier
en $X_s$ (es decir, tales que $D$ es un divisor de Cartier relativo en los
puntos de $X_s$) es abierto en $S$. Además, la definición anterior se formuló
precisamente para que los divisores de Cartier relativos sean estables por todo
cambio de base $S'\longrightarrow S$. Sea $\operatorname{Div}(X/S)$ el
conjunto de divisores relativos sobre $X/S$, y definamos el funtor
contravariante en $S'$ sobre $S$ mediante
\[
\operatorname{Div}_{X/S}(S')=
\operatorname{Div}(X\times_SS'/S').
\]

Supongamos que $X$ es plano y propio sobre $S$. Por la caracterización (ii) de
los divisores de Cartier relativos, $\operatorname{Div}_{X/S}$ puede
considerarse como un subfuntor del funtor $\operatorname{Hilb}_{X/S}$ definido
en [4], IV, y la inclusión
\[
\operatorname{Div}_{X/S}\longrightarrow\operatorname{Hilb}_{X/S}
\]
es «representable por inmersiones abiertas» (cf. [5], IV, 3.13), por las
observaciones precedentes. El teorema principal de existencia de [4], IV, da
entonces:

\src{152}
\result{PROPOSICIÓN}{4.1} Supongamos que $f:X\longrightarrow S$ es proyectivo
y plano. Entonces el funtor $\operatorname{Div}_{X/S}$ es representable; más
precisamente, está representado por un abierto de
$\operatorname{Hilb}_{X/S}$.

Utilizando, para un haz muy amplio fijo $\mathcal O_X(1)$ sobre $X/S$, la
descomposición canónica de $\operatorname{Hilb}_{X/S}$ en los abiertos
$\operatorname{Hilb}^{Q}_{X/S}$ correspondientes a los polinomios de Hilbert
$Q\in\mathbf Q[t]$, se obtiene asimismo la descomposición análoga
\[
\operatorname{Div}_{X/S}=
\coprod_{Q\in\mathbf Q[t]}\operatorname{Div}^{Q}_{X/S}
\]
como unión disjunta de abiertos cuasiproyectivos sobre $S$.

La asignación $D\longmapsto\mathcal L(D)$ proporciona además un homomorfismo
funtorial
\begin{equation}
\operatorname{Div}_{X/S}\longrightarrow\operatorname{Pic}_{X/S}
\tag{+}
\end{equation}
que estudiaremos ahora; bajo condiciones bastante generales resultará
relativamente representable ([5], IV, 3). Partamos de un elemento $\xi$ de
$\operatorname{Pic}_{X/S}(S')$, tomando $S'=S$ para simplificar la notación, y
mostremos que el subfuntor correspondiente de $\operatorname{Div}_{X/S}$ es
representable. Consideremos primero el caso en que $\xi$ viene definido por un
Módulo invertible $\mathcal L$ sobre $X$. Supongamos que $X$ es propio y plano
sobre $S$ y que sus fibras geométricas sobre $S$ son íntegras. Esto implica
también ([2], III, párr. 7) que
$\mathcal O_S\xrightarrow{\sim}f_*(\mathcal O_X)$ y que esta relación sigue
siendo válida tras todo cambio de base $S'\longrightarrow S$. Los divisores de
Cartier relativos $D$ sobre $X/S$ para los que $\mathcal L(D)$ y $\mathcal L$
definen el mismo elemento de
$\operatorname{Pic}(X/S)=\operatorname{Pic}_{X/S}(S)$, o, equivalentemente por
2.4, para los que $\mathcal L(D)$ y $\mathcal L$ son localmente isomorfos sobre
$S$, están en biyección con las secciones del haz cociente
$f_*(\mathcal L)^*/\mathcal O_S^*$. Esta correspondencia es compatible con el
cambio de base. Consideraciones generales de tipo «Künneth» del pasaje citado
(véase también [6]) muestran que la propiedad de ser propio de $X/S$ y la
planitud de $\mathcal L$ sobre $S$ implican la existencia de un Módulo coherente
$\mathcal Q$ sobre $S$, definido salvo isomorfismo único, y de un isomorfismo
de haces
\[
f_*(\mathcal L)\xrightarrow{\sim}
\operatorname{Hom}_{\mathcal O_S}(\mathcal Q,\mathcal O_S).
\]
\footnote{El texto impreso coloca $\mathcal O_X$ como subíndice de
$\operatorname{Hom}$; puesto que $\mathcal Q$ es un Módulo sobre $S$, su dual
se toma sobre $\mathcal O_S$.} La formación de $\mathcal Q$ es compatible con
el cambio de base. Aquí $f_*(\mathcal L)^*$ denota el subhaz de conjuntos de
$f_*(\mathcal L)$ cuyas secciones sobre $U$ son las secciones de $\mathcal L$
sobre $f^{-1}(U)$ que definen divisores de Cartier relativos sobre

\src{153}
$f^{-1}(U)/U$, es decir, que inducen secciones que no son divisores de cero
sobre los $X_s$ para $s\in U$. Como las fibras $X_s$ son íntegras, esto
significa simplemente que las secciones inducidas sobre las fibras $X_s$ no
son idénticamente nulas o, equivalentemente, en términos de homomorfismos
locales $\mathcal Q\longrightarrow\mathcal O_S$, que estos homomorfismos son
sobreyectivos, por el lema de Nakayama. Así, el conjunto de secciones de
$f_*(\mathcal L)^*/\mathcal O_S^*$ está en biyección con el conjunto de
Módulos cociente invertibles de $\mathcal Q$, o, por la definición del fibrado
proyectivo $\mathbf P(\mathcal Q)$ asociado al Módulo coherente $\mathcal Q$
(cf. [5], V, 2), con el conjunto de secciones de $\mathbf P(\mathcal Q)$ sobre
$S$. Esta descripción es compatible con las imágenes inversas y proporciona,
por tanto, el teorema siguiente.

\result{TEOREMA}{4.3} Sea $f:X\longrightarrow S$ un morfismo propio y plano de
fibras geométricas íntegras, donde $S$ es localmente noetheriano, y sea
$\mathcal L$ un Módulo invertible sobre $X$. Para todo $S'$ sobre $S$, sea
$T(S')$ el conjunto de divisores relativos $D$ sobre $X\times_SS'/S'$ tales
que $\mathcal L(D)$ sea isomorfo a
$\mathcal L\otimes_{\mathcal O_S}\mathcal O_{S'}$ localmente sobre $S'$, es decir, tales que
$\mathcal L(D)$ y $\mathcal L\otimes_{\mathcal O_S}\mathcal O_{S'}$ definan
el mismo elemento de $\operatorname{Pic}(X\times_SS'/S')$. Existe entonces un
Módulo coherente $\mathcal Q$ sobre $S$, determinado salvo isomorfismo único,
tal que el funtor $T$ está representado por el fibrado proyectivo
$\mathbf P(\mathcal Q)$.

\result{COROLARIO}{4.4} Si, además, $f$ es proyectivo, el homomorfismo funtorial
$\operatorname{Div}_{X/S}\longrightarrow\operatorname{Pic}_{X/S}$ es
representable por morfismos proyectivos.

En efecto, cuando $X$ admite una sección (respectivamente, localmente una
sección) sobre $S$, el homomorfismo anterior está representado por fibrados
proyectivos (respectivamente, fibrados proyectivos locales), en virtud de 4.3 y
2.1. Cuando $f$ es cuasiproyectivo, se reduce fácilmente al caso anterior
mediante un argumento de descenso que utiliza cuasisecciones finitas y planas
de $X$ localmente sobre $S$.

\result{OBSERVACIÓN}{4.5} Bajo las condiciones de 4.3, el Módulo $\mathcal Q$
sobre $S$ no es, en general, localmente libre, como expresa el hecho de que la
dimensión de las fibras reducidas de $\mathcal Q$, es decir, la dimensión de
$H^0(X_s,\mathcal L_s)$ para $s\in S$ variable, pueda dar saltos. Dado un
Módulo coherente $\mathcal Q$ sobre un preesquema localmente noetheriano $S$,
se comprueba fácilmente que, para un $s\in S$ fijo, $\mathcal Q$ es libre en
$s$ si y solo si $\mathbf P(\mathcal Q)$ es plano sobre $S$ en los puntos
situados sobre $s$; en tal caso es incluso liso sobre $S$ en esos puntos. Aquí,
con $\mathcal Q$ definido en términos de $\mathcal L$ como antes, esto significa
también que la formación de la imagen directa $f_*(\mathcal L)$ «conmuta con el
cambio de base
\src{154}
en un entorno de $s$», o, equivalentemente, que
$f_*(\mathcal L)_s\longrightarrow H^0(X_s,\mathcal L_s)$ es sobreyectivo. Esto
ocurre, por ejemplo, si $H^1(X_s,\mathcal L_s)=0$. Suponiendo que existan los
preesquemas en cuestión, estos criterios se aplican en particular a la situación
universal $\operatorname{Div}_{X/S}\longrightarrow\operatorname{Pic}_{X/S}$ y
dan, respectivamente, una condición necesaria y suficiente y una condición
suficiente para que este morfismo sea liso en un punto dado de
$\operatorname{Div}_{X/S}$.

\fsection{5}{Demostración del teorema principal de existencia}

Bajo las condiciones de 3.1, elijamos un Módulo muy amplio $\mathcal O_X(1)$
sobre $X/S$, y sea $\xi$ el elemento correspondiente de
$\operatorname{Pic}(X/S)$. Pongamos, para abreviar,
\[
P(S')=\operatorname{Pic}(X\times_SS'/S'),
\]
y supongamos primero, para simplificar, que $X/S$ admite una sección. Sea
$P^+(S')$ el subconjunto de $P(S')$ formado por las clases de Módulos
invertibles $\mathcal L$ sobre $X\times_SS'$ tales que
\[
R^if'_*(\mathcal L(n))=0\quad\text{para $i>0$ y todo $n\geq0$},
\qquad
f'_*(\mathcal L(n))\neq0\quad\text{para todo $n\geq0$}.
\]
Estas condiciones son estables por cambio de base y definen, por tanto, un
subfuntor $P^+$ de $P$, manifiestamente estable por traslación por $\xi$.
Utilizando los «teoremas A y B» de Serre ([2], III, 2) y las generalidades de
[5], IV, 5, se ve fácilmente que $P$ es representable si y solo si lo es $P^+$;
en ese caso, $P^+$ está representado por un abierto $U$ del preesquema
$\operatorname{Pic}_{X/S}$ que representa a $P$, y este es la unión creciente
de los abiertos $U-n\xi$.

Pongamos, para abreviar, $D=\operatorname{Div}_{X/S}$, y sea $D^+$ la imagen
inversa de $P^+$ por el morfismo canónico $D\longrightarrow P$. Se tiene así un
morfismo
\begin{equation}
D^+\longrightarrow P^+,
\tag{+}
\end{equation}
y ya se sabe que $D^+$ está representado por un abierto $D^+$ del preesquema
$D=\operatorname{Div}_{X/S}$, cuya existencia garantiza 4.1. Se deduce
fácilmente de 4.3 que (+) está representado por morfismos lisos, proyectivos y
sobreyectivos; más precisamente, por fibrados proyectivos asociados a Módulos
localmente libres y no nulos en ningún punto. En efecto, si $\mathcal L$ sobre
$X\times_SS'$ es como al comienzo de esta sección, $f'_*(\mathcal L)$ es un
Módulo localmente libre no nulo cuya formación conmuta con el cambio de base;
en la notación de 4.3, $\mathcal Q$ es entonces isomorfo al dual de
$f'_*(\mathcal L)$. Deduciremos la representabilidad de $P^+$ del criterio
general [5], IV, 4.7. En el pasaje citado, tomemos para $\mathfrak C$ la clase
de los morfismos de preesquemas fielmente

\src{155}
planos y cuasicompactos, que son epimorfismos efectivos por [4], I B. Acabamos
de verificar la condición (a) del pasaje citado, a saber, que (+) es
representable por morfismos pertenecientes a $\mathfrak C$. También se ha
verificado la condición (b), que dice que el funtor $P^+$ es compatible con el
descenso fielmente plano y cuasicompacto; esto es inmediato. Queda por probar
la condición (c) de dicho pasaje: la relación de equivalencia $R$ en el
preesquema $D^+$ deducida del morfismo (+), representable por $\mathfrak C$, es
$\mathfrak C$-efectiva, es decir, es efectiva y
$D^+\longrightarrow D^+/R$ pertenece a $\mathfrak C$. Obsérvese primero que
los abiertos $D^{+Q}$ de $D^+$ correspondientes a los distintos polinomios de
Hilbert virtuales $Q\in\mathbf Q[t]$ son estables por $R$, pues las fibras de
$R$ son conexas. Basta, por tanto, mostrar, para cada $Q$, que la relación de
equivalencia $R^Q$ inducida en $D^{+Q}$ es $\mathfrak C$-efectiva. Ahora bien,
$D^{+Q}$ es cuasiproyectivo, mientras que la relación de equivalencia $R^Q$ es
proyectiva y plana. Se cumplen, pues, las hipótesis de [4], III, 6.1, que dan el
resultado requerido.

En el caso general en que $X/S$ no admite necesariamente una sección, se reduce
fácilmente al caso anterior por descenso, o se repite la demostración precedente
con la modificación necesaria en la definición de $P^+$.

\result{OBSERVACIONES}{5.1} El método empleado es esencialmente el método de
Matsusaka para la construcción proyectiva de las variedades de Picard. El
resultado de [4], III, invocado para pasar al cociente efectivo, también puede
deducirse fácilmente del teorema de existencia de los esquemas de Hilbert
(véase, por ejemplo, [6]); clásicamente, estos cocientes se construían con
coordenadas de Chow. Obsérvese que la formación del abierto
$\operatorname{Pic}^{+}_{X/S}$ de $\operatorname{Pic}_{X/S}$ y su
descomposición en abiertos $\operatorname{Pic}^{+Q}_{X/S}$ cuasiproyectivos
sobre $S$, según los polinomios de Hilbert de los divisores que definen los
Módulos invertibles considerados, es compatible con el cambio de base, lo que
permite aplicar el argumento de descenso.

\result{OBSERVACIONES}{5.2} Es posible que $\operatorname{Pic}_{X/S}$ exista
siempre que $f:X\longrightarrow S$ sea propio y plano y los homomorfismos
\[
k(s)\longrightarrow H^0(X_s,\mathcal O_{X_s})\qquad(s\in S)
\]
sean isomorfismos; esta última condición significa entonces también que
$\mathcal O_S\xrightarrow{\sim}f_*(\mathcal O_X)$ y que esta relación sigue
siendo válida tras todo cambio de base $S'\longrightarrow S$. Esto es, al
menos, lo que se demuestra para espacios analíticos cuando $f$ es además
proyectivo, mediante un método distinto (el de Chow, salvo error), expuesto en
[5], IX, 3.1.\footnote{La fe de erratas de 1962 indica que la conjetura de existencia
formulada al comienzo de esta observación es falsa, aunque los métodos de
Mumford demuestran un teorema ligeramente más débil.} En ese método, el paso al

\src{156}
cociente por una relación de equivalencia propia y plana en un abierto del
esquema de divisores se sustituye por el paso al cociente por el grupo
proyectivo en el esquema de inmersiones de $X$ en $\mathbf P^r_S$. Este método
puede probablemente adaptarse a los esquemas mediante los resultados de
Mumford sobre el paso al cociente por el grupo proyectivo [8]. Por el momento,
solo existe una demostración escrita cuando $X$ tiene «muchas secciones
locales» sobre $S$, por ejemplo, cuando $X$ es separable sobre un anillo local
completo con cuerpo residual algebraicamente
cerrado.\footnote{El texto impreso omite «con cuerpo residual algebraicamente
cerrado»; la fe de erratas dice que el contraejemplo de Mumford hace indispensable
esta restricción.} En principio, este método tendría mayor alcance, pues también
da la existencia de preesquemas de Picard en casos en los que no son separados
y en los que, por tanto, el primer método necesariamente falla. (Técnicamente,
la dificultad consiste en que, cuando las fibras geométricas de $f$ dejan de
ser íntegras, el funtor considerado en 4.3 no está representado por el propio
fibrado proyectivo $\mathbf P(\mathcal Q)$, sino por un abierto suyo, lo que
conduce al delicado problema de tomar el cociente por una relación de
equivalencia plana pero no propia.)

\result{OBSERVACIÓN}{5.3} Obsérvese que la demostración dada aquí no utiliza ni
una construcción previa de las jacobianas de curvas o familias de curvas, ni la
teoría de variedades abelianas o esquemas abelianos. Difiere, por ello,
esencialmente de exposiciones tradicionales como el libro de Lang [7] o la
exposición de Chevalley [1], que siguen el camino esbozado por A. Weil. Incluso
para las jacobianas de curvas no singulares sobre un cuerpo algebraicamente
cerrado, digamos el cuerpo de los números complejos, la construcción dada aquí
es la única conocida que proporciona las propiedades muy fuertes adoptadas
como definición en la sección 1: esencialmente las de Chevalley, pero permitiendo
«variedades de parámetros» con elementos nilpotentes. Que la construcción de
los esquemas de Picard debe preceder, y no seguir, a la teoría de variedades
abelianas es claro a priori, pues los esquemas de Picard no son, en general,
esquemas abelianos ni se reducen a ellos. Esto ya se ve para curvas singulares
sobre un cuerpo algebraicamente cerrado, donde aparecen las «jacobianas
generalizadas» de Rosenlicht, que no son variedades abelianas. Además, la
teoría de variedades abelianas, y más generalmente de esquemas abelianos, se
simplifica mucho cuando se dispone de una teoría general de esquemas de Picard.
En particular, la teoría de dualidad para esquemas abelianos, especialmente los
resultados de tipo Cartier, se vuelve entonces casi formal (véase, por ejemplo,
[10]).

\src{157}
\result{OBSERVACIONES}{5.4} El «principio de compatibilidad» de Igusa para la
jacobiana de una curva que degenera en una curva singular solo puede
comprenderse adecuadamente como un teorema de existencia del esquema de Picard
de un esquema relativo en curvas $X/S$, no necesariamente liso sobre $S$. Es,
por tanto, un caso particular del teorema principal de existencia 3.1 cuando la
curva especializada es íntegra, es decir, en lenguaje clásico, irreducible de
multiplicidad 1. En la actualidad, el caso de una curva especial reducible,
incluso cuando sus componentes tienen multiplicidad 1, es decir, cuando la
curva especial es separable sobre el cuerpo residual, queda fuera de los
teoremas de existencia conocidos, salvo sobre un anillo de valoración discreta
completo con cuerpo residual algebraicamente cerrado; compárese 5.2. Esta
cuestión de existencia aparecerá ciertamente en una construcción geométrica, en
teoría de esquemas, de las «compactificaciones» de Baily--Satake de los esquemas
de módulos de curvas de género $g$. (Por el momento, esta compactificación solo
se conoce para $g=1$, gracias a los trabajos de Igusa.)

\fsection{6}{Teoremas de existencia relativos}

Esbozaremos aquí algunos casos útiles en los que la existencia de ciertos
esquemas de Picard implica la de otros, lo que permite deducir diversos
teoremas de existencia adicionales del teorema principal 3.1.

\result{PROPOSICIÓN}{6.1} Sea $f:X\longrightarrow S$ un morfismo plano y
proyectivo tal que, en la factorización de Stein $f=f''f'$, el morfismo
$f':X\longrightarrow S'$ es plano, con fibras geométricas íntegras, y satisface
por tanto las hipótesis de 3.1, mientras que el morfismo finito
$f'':S'\longrightarrow S$ es plano. Entonces $\operatorname{Pic}_{X/S}$ existe
y, con la notación introducida en [4], II, C 2, se tiene un isomorfismo canónico
\[
\operatorname{Pic}_{X/S}\xrightarrow{\sim}
\Pi_{S'/S}\operatorname{Pic}_{X/S'}.
\]

Para demostrarlo, se establece primero un isomorfismo de funtores
\[
\operatorname{Pic}_{X/S}\xrightarrow{\sim}
\Pi_{S'/S}\operatorname{Pic}_{X/S'},
\]
se aplica 3.1, que implica que $\operatorname{Pic}_{X/S'}$ es representable, y
se usa después la estructura de $\operatorname{Pic}_{X/S'}$ descrita en la
sección 3, que implica que todo subconjunto finito de una fibra de
$\operatorname{Pic}_{X/S'}$ sobre $S$ está contenido en un abierto afín, para
obtener la existencia de $\Pi_{S'/S}\operatorname{Pic}_{X/S'}$.

\src{158}
Por ejemplo, cuando $X$ es un esquema que es unión disjunta de esquemas $X_i$
sobre $S$ que satisfacen las condiciones de 3.1, el enunciado 6.1 se reduce al
enunciado trivial
\[
\operatorname{Pic}_{X/S}\xrightarrow{\sim}
\prod_i\operatorname{Pic}_{X_i/S}.
\]

\result{COROLARIO}{6.2} Sea $f:X\longrightarrow S$ un morfismo proyectivo y
plano con fibras geométricas localmente íntegras (por ejemplo, un morfismo
proyectivo normal). Entonces $\operatorname{Pic}_{X/S}$ existe.

En este caso, además, $S'$ es un recubrimiento étale de $S$; esto ocurre
siempre que $f$ sea separable, es decir, plano con fibras geométricamente
reducidas. Se comprueba que el teorema de estructura enunciado en la sección 3
para $\operatorname{Pic}_{X/S}$ sigue siendo válido, en virtud de la estructura
análoga de $\operatorname{Pic}_{X/S'}$.

Por otra parte, un procedimiento de descenso proporciona un teorema de
existencia relativo cuyo alcance depende de la solución de las cuestiones de
descenso no plano planteadas en [4], I, A (c). Solo explicitaremos el caso
particular siguiente:

\result{PROPOSICIÓN}{6.3} Sea $f:X\longrightarrow S$ un morfismo propio, y sean
$X_1$ y $X_2$ dos subesquemas de $X$, planos sobre $S$, definidos por dos
Ideales coherentes $\mathcal I_1$ y $\mathcal I_2$ tales que
$\mathcal I_1\cap\mathcal I_2=(0)$ y
$\mathcal O_X/(\mathcal I_1+\mathcal I_2)$ es plano sobre $S$; es decir, el
subesquema de $X$ que es el supremo de $X_1$ y $X_2$ es $X$, mientras que su
ínfimo $Z$ es plano sobre $S$.\footnote{La letra final aparece impresa como
«$X$», pero la condición inmediatamente anterior concierne a la planitud de
$\mathcal O_X/(\mathcal I_1+\mathcal I_2)$ sobre $S$ y, por tanto, define $Z$
como plano sobre $S$.} Supongamos además que, para todo $s\in S$, los
homomorfismos
\[
k(s)\longrightarrow H^0(X_{i,s},\mathcal O_{X_{i,s}}),\qquad i=1,2,
\]
son biyectivos. Entonces el homomorfismo natural de funtores
\[
\operatorname{Pic}_{X/S}\longrightarrow
\operatorname{Pic}_{X_1/S}\times\operatorname{Pic}_{X_2/S}
\]
es representable por morfismos afines. Por consiguiente, si
$\operatorname{Pic}_{X_1/S}$ y $\operatorname{Pic}_{X_2/S}$ existen, también
existe $\operatorname{Pic}_{X/S}$, y el morfismo canónico
\[
\operatorname{Pic}_{X/S}\longrightarrow
\operatorname{Pic}_{X_1/S}\times\operatorname{Pic}_{X_2/S}
\]
es afín.

Por descenso fielmente plano se reduce al caso en que $Z$ admite una sección
sobre $S$, que define así secciones de $X$, $X_1$ y $X_2$ sobre $S$ y permite
eliminar los automorfismos de las estructuras consideradas como se explicó en
2.5. La demostración consiste entonces en observar que los datos de un Módulo
invertible

\src{159}
«rigidificado» $\mathcal L$ sobre $X$ equivalen a los datos de una terna
$(\mathcal L_1,\mathcal L_2,u)$, donde $\mathcal L_i$ es un Módulo invertible
«rigidificado» sobre $X_i$ y $u$ es un isomorfismo de $\mathcal L_1|Z$ con
$\mathcal L_2|Z$ compatible con las rigidificaciones. Solo hay que comprobar
que, para $\mathcal L_1$ y $\mathcal L_2$ fijos, el dato de $u$ se expresa por
una sección de un esquema adecuado sobre $S$ que es afín sobre $S$; esto es
fácil. De 6.3 se deduce inmediatamente:

\result{COROLARIO}{6.4} Sea $X$ un esquema propio y separable sobre un cuerpo
$k$, y sean $X_i$ las componentes irreducibles de $X$. Si existen los
$\operatorname{Pic}_{X_i/k}$, también existe $\operatorname{Pic}_{X/k}$, y el
morfismo canónico
\[
\operatorname{Pic}_{X/k}\longrightarrow\prod_i\operatorname{Pic}_{X_i/k}
\]
es afín.

Combinado con 6.2, esto demuestra, por ejemplo, la existencia de
$\operatorname{Pic}_{X/k}$ siempre que $X$ sea un esquema proyectivo y
separable sobre un cuerpo $k$. Cuando $X$ deja de ser separable sobre $k$,
existe asimismo un resultado de reducción que utiliza un argumento debido a
Oort [11]. El método se aplica también sobre un esquema de base arbitrario, caso
útil, por ejemplo, para demostrar en la exposición siguiente el resultado de
finitud anunciado en 3.3. Para evitar un enunciado excesivamente técnico,
limitémonos a un cuerpo de base:

\result{PROPOSICIÓN}{6.5} Sea $X$ un esquema propio sobre un cuerpo $k$, y sea
$X_0$ un subesquema con el mismo conjunto subyacente, definido por tanto por un
Ideal nilpotente sobre $X$. Entonces el morfismo funtorial
\[
\operatorname{Pic}_{X/k}\longrightarrow\operatorname{Pic}_{X_0/k}
\]
es representable por morfismos afines. En particular, si existe
$\operatorname{Pic}_{X_0/k}$, también existe $\operatorname{Pic}_{X/k}$, y el
morfismo $\operatorname{Pic}_{X/k}\longrightarrow\operatorname{Pic}_{X_0/k}$
es afín.

Combinando esto con 6.4 se concluye inmediatamente:

\result{COROLARIO}{6.6} Sea $X$ un preesquema proyectivo sobre un cuerpo $k$.
Entonces $\operatorname{Pic}_{X/k}$ existe.

\result{OBSERVACIÓN}{6.6} Es sumamente plausible que
$\operatorname{Pic}_{X/k}$ exista para todo esquema propio $X$ sobre un cuerpo
$k$.\footnote{La fe de erratas de 1962 informa de que Murre acababa de responder
afirmativamente a esta cuestión.} Los resultados anteriores reducen esta
cuestión al caso en que $k$ es algebraicamente cerrado y $X$ es íntegro. Por el
lema de Chow, existe entonces un esquema íntegro $X'$, proyectivo sobre $k$, y
un morfismo dominante $g:X'\longrightarrow X$. Bastaría, por tanto, mostrar que
el

\src{160}
morfismo funtorial correspondiente
\[
\operatorname{Pic}_{X/k}\longrightarrow\operatorname{Pic}_{X'/k}
\]
es representable, y sin duda representable por morfismos afines, puesto que ya
se sabe que $\operatorname{Pic}_{X'/k}$ es representable. Esto plantea
cuestiones de descenso no plano que por el momento siguen sin resolverse.
Obsérvese que, si se restringe el funtor $\operatorname{Pic}_{X/k}$ a los
preesquemas reducidos, con $X$ propio e íntegro sobre un cuerpo algebraicamente
cerrado $k$, sí se obtiene un funtor representable, como demostró Chevalley [1]
cuando $X$ es normal y como demostró Seshadri [12] en general mediante un
método de descenso. Sin embargo, con nuestra notación, el esquema construido
por estos autores no es $\operatorname{Pic}_{X/k}$, sino
$(\operatorname{Pic}_{X/k})_{\mathrm{red}}$, el esquema reducido correspondiente
a $\operatorname{Pic}_{X/k}$.

\result{OBSERVACIÓN}{6.7} Más generalmente, sea $f:X'\longrightarrow X$ un
morfismo sobreyectivo de preesquemas propios sobre $k$. Consideraciones de
descenso no plano llevan a conjeturar que
\[
\operatorname{Pic}_{X/k}\longrightarrow\operatorname{Pic}_{X'/k}
\]
es un morfismo afín.\footnote{La fe de erratas de 1962 informa de que esta cuestión
tiene respuesta afirmativa y remite al último párrafo de los Comentarios a la
exposición siguiente.} Esto implicaría, en particular, tras dividir por las
componentes conexas de los elementos neutros, que el homomorfismo
correspondiente de grupos de Néron--Severi es inyectivo módulo torsión. Esto se
puede verificar mediante la teoría de intersecciones cuando $X$ y $X'$ son no
singulares. La respuesta no parece conocida en ningún otro caso.

\result{OBSERVACIÓN}{6.8} Contra lo que podría pensarse, los esquemas de Picard
de esquemas algebraicos con elementos nilpotentes son útiles, e incluso
indispensables, en diversas cuestiones. Así, cuando $X$ es un esquema
proyectivo, digamos liso sobre $k$,\footnote{El texto impreso dice «regular»; la
fe de erratas prescribe «liso sobre $k$».} y $Y$ es una sección hiperplana, hay que
considerar los «entornos infinitesimales» $X_n$ de $Y$ de todos los órdenes y
los esquemas de Picard $\operatorname{Pic}_{X_n/k}$. Cuando $X$ es irreducible
de dimensión $\geq4$ (respectivamente $\geq3$), el morfismo canónico
\[
\operatorname{Pic}_{X/k}\longrightarrow\operatorname{Pic}_{X_n/k}
\qquad(n\text{ grande})
\]
es un isomorfismo (respectivamente, induce un isomorfismo sobre las imágenes
inversas de los subgrupos de torsión de Néron--Severi), resultado que será útil
en el estudio cualitativo de los esquemas de Picard en la exposición siguiente.
Asimismo, considerar los esquemas de Picard de ciertas curvas con elementos
nilpotentes, junto con los teoremas fundamentales de geometría formal [3],
permite verificar, en el caso equicaracterístico, una conjetura de Mumford:
para todo anillo local noetheriano normal completo $A$ de dimensión 2, el grupo
de clases de divisores de $A$ puede considerarse como el conjunto de puntos
racionales sobre $k$ de un

\src{161}
grupo algebraico $G$ sobre el cuerpo residual $k$; además, $G$ queda
determinado canónicamente una vez elegido un cuerpo de coeficientes en $A$.
Cuando $A$ tiene dimensión arbitraria, es plausible que exista un grupo
proalgebraico sobre $k$ que desempeñe el mismo papel que $G$. Cuando
$\operatorname{Spec}(A)$ puede «desingularizarse», dicho grupo se construye como
límite proyectivo de esquemas de Picard de esquemas proyectivos adecuados sobre
$k$ con elementos nilpotentes.

\section*{BIBLIOGRAFÍA}
\addcontentsline{toc}{section}{Bibliografía}

\begin{description}
\item[{[1]}] CHEVALLEY (Claude). -- Sur la théorie de Picard,
\emph{Amer. J. of Math.}, vol. 82, 1960, pp. 435--490.

\item[{[2]}] GROTHENDIECK (A.) y DIEUDONNÉ (J.). -- Éléments de géométrie
algébrique, Chapitre I et suivants. -- París, Presses universitaires de France,
1960--1961 (Institut des Hautes Études Scientifiques, Publications
mathématiques 4, 8, \ldots).

\item[{[3]}] GROTHENDIECK (Alexander). -- Géométrie formelle et géométrie
algébrique, \emph{Séminaire Bourbaki}, vol. 11, 1958/59,
n.\textsuperscript{o} 182, 28 pp.

\item[{[4]}] GROTHENDIECK (Alexander). -- Techniques de descente et théorèmes
d'existence en géométrie algébrique, I--IV, \emph{Séminaire Bourbaki}, vol. 12,
1959/60, n.\textsuperscript{o} 190, 29 pp. y n.\textsuperscript{o} 195, 22 pp.;
vol. 13, 1960/61, n.\textsuperscript{o} 212, 20 pp. y
n.\textsuperscript{o} 221, 28 pp.

\item[{[5]}] GROTHENDIECK (Alexander). -- Techniques de construction en
géométrie analytique, I--X, \emph{Séminaire Cartan}, vol. 13, 1960/61:
Déformation des structures analytiques complexes,
n.\textsuperscript{o} 7--17.

\item[{[6]}] GROTHENDIECK (Alexander). -- Séminaire de géométrie algébrique,
editado por Lichtenbaum, Harvard University, 1961 (de próxima aparición).

\item[{[7]}] LANG (Serge). -- \emph{Abelian varieties}. -- Nueva York,
Interscience Publishers, 1959 (Interscience Tracts in pure and applied
Mathematics, 7).

\item[{[8]}] MUMFORD (D.). -- An elementary theorem in geometric invariant
theory, \emph{Bull. Amer. math. Soc.}, vol. 67, 1961, pp. 483--486.

\item[{[9]}] MUMFORD (D.). -- The topology of normal singularities of an
algebraic surface and a criterion for simplicity. -- París, Presses
universitaires de France, 1961 (Institut des Hautes Études Scientifiques,
Publications mathématiques, 9).

\item[{[10]}] MUMFORD (D.) y TATE (J.). -- Séminaire de géométrie algébrique,
Harvard University, Spring Term 1962 (de próxima aparición).

\item[{[11]}] OORT (F.). -- Sur le schéma de Picard,
\emph{Bull. Soc. math. France}, vol. 90, 1962 (de próxima aparición).

\item[{[12]}] SESHADRI (C. S.). -- Tesis, de próxima aparición en
\emph{Annali di Matematica}.
\end{description}
